\section{Current assumptions}\label{sec:assumptions}
The simplifications the derivation and the simulator rely on. Each is expected to be relaxed
later, unless stated otherwise.
\begin{itemize}
    \item \textbf{Planar flight, no crosswind.} Motion is confined to the vertical plane fixed by
    the launch azimuth $\chi$; the out-of-plane wind component is computed and discarded. Wind is
    horizontal only: updrafts and downdrafts are not modelled. The yaw rate is
    $\omega_z\equiv0$, which is what makes the roll freedom of \autoref{sec:roll} decouple
    exactly rather than approximately.
    \item \textbf{Gravity direction.} Parallel by default, $\hat{z}_g\equiv\hat{z}_0$. The radial
    and ellipsoidal direction models are available; once one is enabled the ground is curved and
    the true altitude $h_{true}$ (\autoref{eq:h_true}) replaces $h=-z$, with impact at
    $h_{true}=0$.
    \item \textbf{Earth rotation, in part.} The centrifugal effect on the magnitude of $g$ and the
    in-plane Coriolis term are included. The out-of-plane Coriolis component and the residual
    centrifugal term are dropped, and $\vec{\Omega}$ is frozen at the launch site's latitude and
    heading.
    \item \textbf{Moment arms.} Mass, centre of gravity and inertia vary with the motor model. The
    centre of pressure $l_p$ responds to the flight state only under the table aerodynamic model;
    the thrust point $l_t$ is fixed.
    \item \textbf{Thrust axis.} The thrust angle $\theta_t$ and the radial offsets $r_g,r_t$ are
    constant inputs, zero by default. A commanded gimbal is not modelled.
    \item \textbf{Rail, recovery, air brake.} The rail has no friction and no tip-off. The canopy
    opens instantly and the attitude is frozen under it. The air brake affects the axial force
    only, through a memoryless law; a law with memory is expressible in the sampled mode but not
    written.
    \item \textbf{Turbulence.} The Dryden spectrum is driven at a fixed reference airspeed, not the
    rocket's own, and has no vertical component. It is off by default.
    \item \textbf{Roll.} The roll freedom of \autoref{sec:roll} is exact in the pitch plane, but
    the yaw torques that roll and the wheel produce --- \autoref{eq:roll_gyro_pitch}'s
    $(I_x-I)\omega_z\omega_x$ and \autoref{eq:roll_gyro_wheel}'s $\omega_rh_w$ --- cannot be
    represented at all and are dropped with $\omega_z$. The aerodynamic coefficients do not depend
    on the roll angle, so fin--body interference at incidence, roll lock-in and the Magnus side
    force are absent; the wheel is balanced and centred on the axis, so spinning it moves neither
    the centre of gravity nor the pitch inertia. Off by default.
    \item \textbf{Three-dimensional motion.} \autoref{sec:rocket3d} lifts the two items above ---
    planar flight and the planar roll freedom --- but it is a formulation of the rotational
    dynamics, not yet a flyable mode, and it carries assumptions of its own. The inertia tensor
    (\autoref{eq:inertia_matrix}) is \emph{constant}: the 3D formulation is written for a rigid body
    of fixed mass distribution, so the variable-mass inertia $I(t)$ of \autoref{eq:mass_inertia} and
    the term $\dot{\mathbf{I}}\vec\omega$ it would add are both absent. The moment vector
    (\autoref{eq:moment_3d}) supplies the pitching component only: there is no yaw moment from
    sideslip and no roll moment from fin cant, so a 3D run with a symmetric tensor reproduces the
    planar pitch history and adds nothing. A reaction wheel, if present, is taken as locked (its
    inertia inside $I_{xx}$, its momentum outside $\vec{H}$). The rail constraint, the canopy phase
    and the flight-phase events are all written against the planar state
    (\autoref{eq:state_vector}) and have no 3D counterpart, so a 3D scenario is integrated over a
    fixed time span rather than flown to the ground. Off by default; the planar model remains the
    one every result in this document is produced with.
\end{itemize}
